% file: 拓展.tex
% 作为主文档 \input{../../拓展/拓展} 的子文件，不要写 \documentclass、\begin{document}、\end{document}
% 由同目录 拓展.md 生成（2026-10-20）

\section{拓展}
\noindent\textit{本部分为补充篇：伴随矩阵的秩、降阶公式，以及两处待补的残稿。}

\subsection{伴随矩阵的秩与原矩阵的秩的关系}
设 $A$ 为 $n$ 阶方阵，$A^{*}$ 为其伴随矩阵，则
\[
r(A^{*})=
\begin{cases}
n, & r(A)=n,\\[1mm]
1, & r(A)=n-1,\\[1mm]
0, & r(A)\le n-2.
\end{cases}
\]

\noindent\textit{（以下直觉解释原图没有，2026-10-20 按标准结论补：$r(A)=n$ 时 $A$ 可逆、$A^{*}=|A|A^{-1}$ 也可逆；
$r(A)=n-1$ 时由 $AA^{*}=|A|E=O$ 得 $r(A)+r(A^{*})\le n$，又存在 $n-1$ 阶非零子式使 $r(A^{*})\ge1$；
$r(A)\le n-2$ 时所有 $n-1$ 阶子式为零，$A^{*}=O$。）}

\subsection{降阶公式}
设 $A$ 是 $m\times n$ 矩阵，$B$ 是 $n\times m$ 矩阵，$m>n$，$\lambda$ 为任意数，则
\[
|\lambda E_m-AB|=\lambda^{m-n}\,|\lambda E_n-BA|.
\]

\subsection{待补}
原稿中「生成的向量空间」只有半句定义，「相似的本质」只有标题，均未写出内容，按原样保留待补。
